% !TEX root = ../main.tex

\section{Introduction}

Automated text evaluation is increasingly moving from simple similarity scoring toward the use of large language models as general-purpose evaluators. Instead of measuring only lexical overlap or semantic proximity to a reference text, LLM-based evaluators can assign scores to open-ended artefacts across flexible, task-specific criteria. This makes them attractive not only for evaluating natural language generation systems, but also for settings in which text must be quantified as evidence about a broader construct: the quality of an answer, the usefulness of a description, or the attractiveness of a marketplace offer. In these settings, the evaluator is not merely asked to compare two strings. It is asked to produce a measurement that supports interpretation, diagnosis, and downstream decision-making.

A central difficulty is that such constructs are rarely unidimensional. A text can be fluent but uninformative, persuasive but misleading, detailed but irrelevant, or attractive to one audience while imposing substantial participation costs on another. Recent work has therefore moved toward multi-dimensional LLM evaluation, where the target is decomposed into criterion-level judgments. This development is valuable because it makes evaluation more diagnostic: instead of producing only an overall score, the evaluator can indicate which aspects of an artefact are strong or weak. However, decomposition alone does not solve the measurement problem. If criteria are selected pragmatically and then recomposed into a single score without an explicit account of how they define the evaluation target, the resulting measurement may be predictive or rating-aligned while still being theoretically underspecified.

This thesis argues that artefact-level LLM evaluation should therefore be treated as a construct-specification problem rather than merely an aggregation problem. The central issue is not only how to elicit reliable criterion scores from an LLM, but whether those criteria validly define the construct that the final score is interpreted as measuring. This distinction matters because many existing LLM-as-judge methods implicitly treat human holistic ratings as the object of measurement. Alignment with such ratings is useful validity evidence, but it does not by itself establish that the evaluator measures the intended construct. In diagnostic applications, especially when the goal is to understand why an artefact performs well or poorly, the object of measurement must be specified independently of the rating procedure.

To formalize this approach, artefact-level scoring is specified as a formative measurement model. In a formative model, the observed components are not interchangeable indicators of an underlying latent trait; rather, they jointly define the construct. For example, the attractiveness of a creator-marketplace deal may be formed by the value of the reward, the costs imposed on the creator, and the quality with which the offer is communicated. None of these dimensions is merely a noisy reflection of an underlying latent attractiveness variable. Instead, changing one dimension changes the construct itself. This makes formative specification particularly suitable for diagnostic artefact-level evaluation, because it preserves criterion-level interpretability while allowing the dimensions to be recomposed into an overall score.

The empirical question of this thesis is whether theoretically grounded formative decomposition produces more valid and interpretable artefact-level measurements than holistic or weakly specified alternatives. This question is studied across two datasets. First, FeedbackQA is used to examine convergent validity: whether LLM-derived scores align with existing human quality judgments for academic question-answer feedback. Second, a proprietary dataset from Barter, a creator marketplace, is used to examine predictive validity: whether LLM-derived measurements of deal attractiveness explain and predict creator applications after controlling for structural deal characteristics. The Barter setting is especially useful because the target construct is not simply text quality in the abstract. It is a domain-specific construct with observable behavioral consequences: creators either apply to collaborate or they do not.

Across both settings, the thesis compares holistic, weakly specified multi-dimensional, and theoretically grounded formative prompting conditions. Rather than relying only on categorical model outputs, it extracts expected values from restricted rating-token logits and uses entropy as an additional diagnostic of rating uncertainty. The resulting scores are evaluated using validity-oriented statistical models: convergent validity for FeedbackQA and predictive validity for Barter Deals. This design separates three questions that are often conflated in LLM evaluation: whether an LLM can assign plausible scores, whether the selected criteria define the intended construct, and whether recomposing those criteria improves downstream validity.

The contribution of this thesis is therefore both methodological and conceptual. Methodologically, it provides an evaluation workflow for using LLMs as artefact-level measurement instruments: define the construct, specify its formative components, elicit criterion-level scores separately, and validate the recomposed measurement against theoretically relevant evidence. Conceptually, it shifts the focus from whether LLM judgments resemble human judgments to whether LLM-based measurements support justified claims about the artefacts being evaluated. This shift is necessary if LLM evaluators are to be used not only as scalable annotators, but as interpretable measurement instruments for open-ended text.


% \subsection{Industrial Partner: Barter}

% The current study is motivated by the operational challenges faced by the industrial partner, Barter. The company aims to leverage Generative AI (GenAI) to streamline the onboarding experience, specifically addressing the friction new partners experience during the sign-up flow. Currently, users must manually write marketing copy such as company descriptions and deal offers before they can make use of the platform. This creates a significant barrier to entry, resulting in user attrition.

% While GenAI offers a scalable solution to automate this drafting process, its implementation in business contexts is hindered by the inherent variability of output quality. As LLMs are probabilistic in nature, some degree of inconsistency is unavoidable. Although well-engineered prompts might yield high-quality output in the majority of cases, the accumulation of low-quality output (such as hallucinations and misalignment) becomes a significant operational problem at high volumes. 

% Two fundamental challenges lie at the core of this problem. First is the challenge of identification: without a reliable mechanism to flag low-quality output, we are unable to quantify the severity of the problem, or make the problematic output actionable. The second is the challenge of diagnosis: improving the output requires an evaluation to identify the reasons why the quality is insufficient. For GenAI to function as an automated solution, it must not only detect the errors but also provide the feedback to resolve them.

% Therefore, although Barter's problem statement concerns generation, the solution to the problem begins with accurate validation. Currently, Barter operates in a similar fashion, where employees use GenAI to generate content which is manually evaluated and refined. However, such a human-in-the-loop (HITL) approach is costly and unscalable. To address this, this thesis explores whether small-scale open-source LLMs can be used for automated "AI-in-the-loop" (AITL) evaluation. We introduce a method inspired by the latent variable models of Psychometrics to provide the granular evaluation required for the identification and diagnosis of output quality. 
